% Adapted source-splice proof for the coefficient-six parent.
\section{Exact source splice and its physical realization}

At a successful MP comparison, the source has $T=H=X=Y=0$, $D=b$, and
arbitrary stack $W=w$. The original subtree is
\[
 \texttt{node163}=(\texttt{node6},\texttt{node55}),\qquad
 \texttt{node55}=(\texttt{node49},\texttt{node54}).
\]
The first child of node55 pushes the current zero; node54 copies $D$ to
$T$ while preserving $D$. PA194N6 replaces node163 by
$(\texttt{node6},\texttt{node54})$.
The deleted source interval is $[3756,3775)$ in the original zero-based
numbering. It contains 19 instructions and has exactly one exterior
entry, the zero branch of instruction 3750. Every remaining instruction,
including every successor, equals the old instruction under the explicit
index relocation in \texttt{construction/build.py}, with that entry redirected
to the retained copy loop. Both failed MP and every other command retain
their source effects.

The retained copy loop, in local labels, is
\[
\begin{array}{c|l}
0&\mathrm{DEC}\ D\quad(0\mapsto4,\ >0\mapsto1)\\
1&\mathrm{INC}\ T;\ \mathrm{goto}\ 2\\
2&\mathrm{INC}\ X;\ \mathrm{goto}\ 3\\
3&\mathrm{goto}\ 0\\
4&\mathrm{DEC}\ X\quad(0\mapsto7,\ >0\mapsto5)\\
5&\mathrm{INC}\ D;\ \mathrm{goto}\ 6\\
6&\mathrm{goto}\ 4
\end{array}
\]
Label seven is the exit. After $j\le b$ iterations of the first loop,
$(T,D,X)=(j,b-j,j)$. After its final zero test and $j\le b$ iterations
of the second, $(T,D,X)=(b,j,b-j)$. Hence the exit is
$(T,D,X)=(b,b,0)$, with all other registers unchanged, in exactly
$7b+2$ primitive source steps. These two elementary induction invariants
prove the retained loop's contract over all naturals.

The literal physical table is rebuilt from the pinned 202-control
radix-six parent and its eight simultaneous folds; the grammar and
resolved source are bound independently. The changed native transition enters the already existing reader
for node54 directly. All 4,877 arbitrary-register primitive outcomes pass
(the maximum symbolic-checker step count is 499); all 32 local fold
contexts pass simultaneously; the 23 arithmetic-engine controls retain
every original row. The complete blank-tape initialization reaches native
control 84 at head $-2$, with the entire clean frame and 26 spare blocks,
after exactly 58,320 literal transitions. The startup equality includes
every nonzero tape cell.

The generic register-block and unary-scan arguments are elementary
inductions over coded finite lists and natural lengths. Their trace
witnesses are primitive recursively constructible. The certificates instantiate those induction arguments at every
primitive source case, giving PA-provable positive finite source simulation.
The independent checker rederives four register-list lemmas from the
literal rows and checks all 4,877 outcomes, using 1,115,525 symbolic
operations; every certified duration has positive constant term, at least
three. These counts concern certificate execution, not a uniform physical
running time. Conversely, given a physical first halt at time $T$, iterate the totalized
source step function for $T+1$ instructions. If this bounded iteration
finds a source halt, reflection is established. Otherwise its positive
nonhalting simulations, preceded by startup, produce a deterministic
physical trace beyond that first halt, a contradiction. This proves
reflection by a bounded construction.
Combining source calibration and physical simulation yields the ordinary
mathematical conclusion
\[
 \mathrm{PA}\vdash
 \bigl(\neg\operatorname{Halt}(\mathrm{PA194N6})
                 \leftrightarrow\operatorname{Con}(\mathrm{PA})\bigr).
\]
The finite and symbolic executables support this argument and bind it to
the exact table. They do not replace the explicitly identified inherited
logical proof transformations, and no newly compiled end-to-end trusted
kernel endpoint is claimed.
